\documentclass[12pt,a4paper]{article}
\usepackage{preamble}

\hypersetup{pdftitle={Шаблон конспекта}}

\begin{document}
% !TEX root = template.tex
% Заполните поля ниже --- шапка соберётся сама.
\renewcommand{\notesuniversity}{Санкт-Петербургский государственный университет}
\renewcommand{\notesfaculty}{Факультет математики и механики}
\renewcommand{\notestitle}{Название конспекта}
\renewcommand{\notessubtitle}{Лекции}
\renewcommand{\notesauthor}{}
\renewcommand{\notesdate}{осенний семестр 2026}
% колонтитул слева: курс · тип/номер · дата
\renewcommand{\notesheader}{Название конспекта \(\cdot\) Лекции \(\cdot\) осенний семестр 2026}

\notesmaketitle


\section{Занятие}

Краткий пример того, как набирать конспект: определения и теоремы ---
обычным \texttt{amsthm}, обычное замечание --- жирным заголовком,
напоминание и отступление --- с чёрной полоской слева.
Формулы нумеруются внутри секции.

\dfn{Гладкая функция}{
Пусть \(\Omega\subset\R^n\) --- область. Говорят, что
\(u\in C^\infty(\Omega)\), если все частные производные \(u\)
существуют и непрерывны на \(\Omega\).
}

\thm{Формула Ньютона--Лейбница}{
Если \(f\in C^1[a,b]\), то
\begin{equation}
    \label{eq:newton-leibniz}
    \int_a^b f'(x)\,dx = f(b)-f(a).
\end{equation}
}

\prf{
Достаточно рассмотреть первообразную \(F(x)=f(x)\) и воспользоваться
определением определённого интеграла.
}

Формулу Ньютона--Лейбница см.~\eqref{eq:newton-leibniz}.

\lem{Оценка}{
Для любых \(x,y\in\R\) выполнено \(\abs{x+y}\le\abs{x}+\abs{y}\).
}

\stat{Следствие треугольника}{
\(\abs{\abs{x}-\abs{y}}\le\abs{x-y}\).
}

\coroll{}{
Неравенство треугольника сохраняется при любом конечном числе слагаемых.
}

\subsection{Вычисления}

Разложение можно писать в \texttt{align} --- номер ставится автоматически:
\begin{align}
    (x+y)^3
    &= (x+y)^2(x+y) \\
    &= (x^2+2xy+y^2)(x+y) \\
    &= x^3+3x^2y+3xy^2+y^3.
\end{align}

\ex{Производная степени}{
\(\dfrac{d}{dx}x^n = n x^{n-1}\) при \(n\in\N\).
}

\exr{Дифференцирование произведения}{
Найдите \(\dfrac{d}{dx}(x\sin x)\).
}

\rem{Напоминание}{
Непрерывность в \(\mathcal{D}(\Omega)\) равносильна секвенциальной
непрерывности: достаточно проверять пределы последовательностей.
}

\note{Обычное замечание --- только жирный заголовок, без вертикальной полоски.}

\section{Графики и диаграммы}

Простой график через \texttt{pgfplots}:
\begin{center}
\begin{tikzpicture}
\begin{axis}[
    width=0.72\textwidth,
    height=5.4cm,
    axis lines=middle,
    xlabel=\(x\),
    ylabel=\(y\),
    xtick={-1,1},
    ytick={-1,1},
    xmin=-2.2, xmax=2.2,
    ymin=-1.4, ymax=1.4,
    every axis plot/.append style={thick, black},
    tick label style={font=\small},
    label style={font=\small},
]
\addplot[domain=-2:2, samples=120] {sin(deg(x))};
\end{axis}
\end{tikzpicture}
\end{center}

Коммутативная диаграмма через \texttt{tikz-cd}:
\begin{center}
\begin{tikzcd}
    A \arrow[r, "f"] \arrow[d, "h"']
        & B \arrow[d, "g"] \\
    C \arrow[r, "k"']
        & D
\end{tikzcd}
\end{center}

Набросок в чистом \texttt{tikz}:
\begin{center}
\begin{tikzpicture}[x=1.1cm, y=1.1cm]
    \draw[-{Stealth}, gray!65] (-0.2,0) -- (3.6,0) node[right] {\(x\)};
    \draw[-{Stealth}, gray!65] (0,-0.2) -- (0,2.3) node[above] {\(y\)};
    \draw[thick] (0.2,0.35) -- (1.4,1.55) -- (3.1,0.7);
    \fill (1.4,1.55) circle (1.5pt);
\end{tikzpicture}
\end{center}

\end{document}
